% =====================================================================
%  ΚΕΦΑΛΑΙΟ 2 — ΘΕΩΡΗΤΙΚΟ ΥΠΟΒΑΘΡΟ
%  Κανόνες που ισχύουν εδώ (Κεφ. 1 \S1.6):
%   - πρωτογενής πηγή για κάθε φράγμα και κάθε ορισμό
%   - καμία αυτούσια παράθεση από άρθρο που δεν βρίσκεται στο corpus/
%   - το πλαίσιο εγκυρότητας ΔΕΝ μπαίνει εδώ. Ανήκει στο κεφάλαιο της
%     μεθοδολογίας, όπως υπόσχεται το \S1.4.
%
%  Η ΑΓΓΛΙΚΗ ΕΙΝΑΙ ΤΟ ΠΡΩΤΟΤΥΠΟ. Αυτό το αρχείο ακολουθεί το
%  content/en/ch2.tex, μπλοκ προς μπλοκ.
% =====================================================================

%%blk background.head from=d5c676003f4b
\chapter{Θεωρητικό υπόβαθρο: μοντέλα, φράγματα και κριτήρια ορθότητας}
\label{ch:background}

%%blk background.intro.01 from=23b735575924
Το Κεφάλαιο~\ref{ch:intro} έδειξε ότι κάθε πρωτόκολλο συναίνεσης είναι
μια θέση σε έναν χώρο σχεδίασης. Αυτό το κεφάλαιο ορίζει τους άξονες του
χώρου. Χωρίς αυτούς, οι συγκρίσεις των επόμενων κεφαλαίων δεν έχουν
μονάδες.

%%blk background.intro.02 from=8545cc0077b9
Η σειρά είναι σκόπιμη. Πρώτα έρχεται το περιβάλλον στο οποίο τρέχει ένα
πρωτόκολλο, δηλαδή ο συγχρονισμός και τα σφάλματα. Έπειτα έρχεται το
πρόβλημα που λύνει. Έπειτα έρχονται τα φράγματα που κανένα πρωτόκολλο
δεν παραβιάζει. Τέλος έρχονται τα κριτήρια με τα οποία κρίνεται η
ορθότητα και η επίδοσή του.

%%blk background.synchrony.head from=751c57e4c1a0
\section{Δύο άξονες συγχρονισμού, όχι ένας}
\label{sec:synchrony}

%%blk background.synchrony.01 from=431e54fae8e8
Η βιβλιογραφία συχνά μιλά για «σύγχρονα» και «ασύγχρονα» συστήματα σαν
να υπάρχει ένας άξονας. Οι Dwork, Lynch και Stockmeyer χωρίζουν δύο
\cite{dwork1988consensus}:

%%blk background.synchrony.axes from=0339863ff5da
\begin{enumerate}[label=\textbf{(\alph*)},leftmargin=2.4em]
  \item \textbf{Η επικοινωνία.} Υπάρχει άνω φράγμα $\Delta$ στον χρόνο
  παράδοσης ενός μηνύματος;
  \item \textbf{Οι επεξεργαστές.} Υπάρχει άνω φράγμα $\Phi$ στον λόγο με
  τον οποίο το ρολόι ενός κόμβου τρέχει γρηγορότερα από ενός άλλου;
\end{enumerate}

%%blk background.synchrony.02 from=060fae9782ef
Ένα σύστημα είναι \emph{σύγχρονο} όταν και τα δύο φράγματα υπάρχουν και
είναι γνωστά εκ των προτέρων. Είναι \emph{ασύγχρονο} όταν \emph{δεν
υπάρχει κανένα} από τα δύο. Οι δύο άξονες είναι ανεξάρτητοι, άρα
ανάμεσά τους υπάρχουν και μεικτά μοντέλα, όπου η μία διάσταση είναι
σύγχρονη και η άλλη όχι. Ο διαχωρισμός δεν είναι σχολαστικισμός.
Ακριβώς γι' αυτόν τον λόγο οι ίδιοι συγγραφείς δίνουν πίνακα
ανθεκτικότητας με πέντε στήλες αντί για δύο, και τρεις από αυτές
αναπαράγει ο Πίνακας~\ref{tab:dls}.

%%blk background.synchrony.sub.partial from=ac00c81ae549
\subsection{Οι δύο μορφές μερικής συγχρονίας}

%%blk background.synchrony.03 from=60a8c8ba8645
Ανάμεσα στα δύο άκρα βρίσκεται η \emph{μερική συγχρονία}. Το ίδιο άρθρο
την ορίζει με δύο τρόπους, και η διαφορά τους έχει σημασία για την
εργασία.

%%blk background.synchrony.04 from=1ab78664fa06
Στην \textbf{πρώτη μορφή} το φράγμα $\Delta$ υπάρχει, αλλά δεν το
ξέρουμε. Ζητάμε πρωτόκολλο που λειτουργεί σωστά όποια κι αν είναι η
τιμή, δηλαδή πρωτόκολλο που δεν έχει το $\Delta$ ενσωματωμένο.

%%blk background.synchrony.05 from=4da263bdf928
Στη \textbf{δεύτερη μορφή} ξέρουμε το $\Delta$, αλλά το σύστημα
μηνυμάτων δεν το σέβεται πάντοτε. Οι συγγραφείς προσθέτουν έναν
περιορισμό. Σε κάθε εκτέλεση υπάρχει μια χρονική στιγμή που οι κόμβοι
δεν γνωρίζουν. Από εκείνη τη στιγμή και μετά, το σύστημα μηνυμάτων
σέβεται το $\Delta$. Την ονομάζουν \en{global stabilization time} και τη
συντομογραφούν GST.

%%blk background.synchrony.06 from=ac2470f8b4e9
Οι ίδιες δύο μορφές ορίζονται και για το $\Phi$. Γι' αυτόν τον λόγο ο
πίνακας του άρθρου έχει χωριστή στήλη για μερικώς σύγχρονους
επεξεργαστές με σύγχρονη επικοινωνία. Η εργασία χρησιμοποιεί τη δεύτερη
μορφή και για τους δύο άξονες, όπως και τα πρωτόκολλα που αποτιμά.

%%blk background.synchrony.fig.synchrony from=dc51e1fb5fbf
\dgfigure{fig:synchrony}%
{Τα τρία μοντέλα συγχρονισμού, με μία κάθετη γραμμή ανά μήνυμα.}%
{Τα τρία μοντέλα συγχρονισμού, με μία κάθετη γραμμή ανά μήνυμα. Στο
σύγχρονο μοντέλο κάθε μήνυμα φτάνει εντός γνωστού $\Delta$. Στο
ασύγχρονο δεν υπάρχει φράγμα. Στο μερικώς σύγχρονο υπάρχει άγνωστη
στιγμή GST, μετά την οποία ισχύει το $\Delta$ \cite{dwork1988consensus}.
Η GST είναι \textbf{μοντέλο ανάλυσης} και όχι μηχανισμός. Κανένας
κώδικας δεν την υπολογίζει. Ο όρος ανήκει στο ίδιο το άρθρο του 1988 και
δεν είναι μεταγενέστερος.}%
{% fig-synchrony — the three timing models, one vertical stroke per message.
% Structure figure: dgflow. dgCFT marks a message that arrives inside a known
% bound, dgBFT one that does not — the standard's own key/bad pairing.
% The axis is 7.5 cm long so that the whole figure, row labels included, comes
% out at \dgwidth like every other figure; the message positions are the
% original ones scaled by 7.5/9.6.
\begin{dgflow}
  % --- synchronous -------------------------------------------------------
  \draw[dgaxis] (0,3.0) -- (7.5,3.0);
  \node[anchor=east, inner sep=4pt] at (-0.1,3.0) {\fl{sync.model.synchronous}};
  \node[dgnote, anchor=north west] at (0,2.84) {\fl{sync.note.synchronous}};
  \foreach \x in {0.781,1.875,2.969,4.063,5.156,6.250}{
    \draw[dgarrow,-,draw=dgCFT] (\x,2.88) -- (\x,3.12);}

  % --- asynchronous ------------------------------------------------------
  \draw[dgaxis] (0,1.6) -- (7.5,1.6);
  \node[anchor=east, inner sep=4pt] at (-0.1,1.6) {\fl{sync.model.asynchronous}};
  \node[dgnote, anchor=north west] at (0,1.44) {\fl{sync.note.asynchronous}};
  \foreach \x in {0.469,0.859,2.344,5.078,5.313,6.719}{
    \draw[dgarrow,-,draw=dgBFT] (\x,1.48) -- (\x,1.72);}

  % --- partially synchronous --------------------------------------------
  \draw[dgaxis] (0,0.2) -- (7.5,0.2);
  \node[anchor=east, inner sep=4pt] at (-0.1,0.2) {\fl{sync.model.partial}};
  \foreach \x in {0.469,1.172,2.031,2.891}{
    \draw[dgarrow,-,draw=dgBFT] (\x,0.08) -- (\x,0.32);}
  \foreach \x in {4.063,4.922,5.781,6.641}{
    \draw[dgarrow,-,draw=dgCFT] (\x,0.08) -- (\x,0.32);}
  % The dashed line stops above the note: in the original it ran through it.
  \draw[dgaxis,-,dashed] (3.594,-0.45) -- (3.594,0.85);
  \node[dgtag, fill=dgPaper, text=dgAux, anchor=south] at (3.594,0.85) {\textbf{GST}};
  \node[dgtag, fill=dgPaper, text=dgBFT, anchor=north] at (1.62,-0.05)
       {\fl{sync.liveness.none}};
  \node[dgtag, fill=dgPaper, text=dgCFT, anchor=north] at (5.60,-0.05)
       {$\le\Delta$, \fl{sync.liveness.guaranteed}};
  \node[dgnote, anchor=north west] at (-1.2,-0.75) {\fl{sync.note.safety}};
\end{dgflow}
}

%%blk background.synchrony.07 from=31ff8b953ef1
Η δεύτερη μορφή είναι εκείνη που υιοθετούν τα βυζαντινά πρωτόκολλα της
εργασίας \cite{yin2019hotstuff} και \cite{malkhi2023hotstuff2}\prep. Η
πρακτική συνέπεια φαίνεται στο Σχήμα~\ref{fig:synchrony}. Η ασφάλεια
ισχύει σε ολόκληρο τον άξονα. Η ζωτικότητα εξασφαλίζεται μόνο δεξιά της
GST.

%%blk background.faults.head from=95886099d973
\section{Μοντέλα σφαλμάτων}
\label{sec:faults}

%%blk background.faults.01 from=b1275f514326
Ο δεύτερος άξονας ρωτά τι επιτρέπεται να κάνει ένας χαλασμένος κόμβος. Η
εργασία χρησιμοποιεί τέσσερις βαθμίδες, με αυξανόμενη κακοήθεια.

%%blk background.faults.02 from=5dc6adc7e0a5
Μια σύμβαση συμβόλων, μία φορά για όλο το κεφάλαιο. Τα $t$ και $f$
δηλώνουν το \emph{ίδιο} μέγεθος, το ανώτατο πλήθος ελαττωματικών κόμβων.
Το κεφάλαιο διατηρεί το σύμβολο κάθε πρωτογενούς πηγής, ώστε οι τιμές να
είναι αναγνωρίσιμες σε αυτήν.

%%blk background.faults.list from=ea920960c543
\begin{description}[leftmargin=2.2em,style=nextline,font=\normalfont\bfseries]
  \item[Κατάρρευση (\en{crash}).] Ο κόμβος σταματά και δεν ξαναμιλά. Το
  σφάλμα είναι \emph{απουσία} πληροφορίας. Το ειδικότερο μοντέλο
  \en{fail-stop} προσθέτει ότι οι υπόλοιποι κόμβοι αντιλαμβάνονται
  \emph{αξιόπιστα} τη διακοπή. Η διάκριση επανέρχεται στην
  Ενότητα~\ref{sec:smr}, όπου το φράγμα $t{+}1$ του Schneider
  προϋποθέτει \en{fail-stop} και όχι σκέτο \en{crash}.
  \item[Παραλείψεις (\en{omission}).] Ο κόμβος χάνει μηνύματα, αλλά όσα
  στέλνει είναι σωστά.
  \item[Βυζαντινό με αυθεντικοποίηση.] Ο κόμβος συμπεριφέρεται αυθαίρετα,
  αλλά δεν μπορεί να πλαστογραφήσει την υπογραφή άλλου κόμβου.
  \item[Βυζαντινό χωρίς αυθεντικοποίηση.] Καμία εγγύηση. Ο κόμβος λέει
  ό,τι θέλει σε όποιον θέλει, και συνεννοείται με τους υπόλοιπους
  ελαττωματικούς \cite{lamport1982byzantine}.
\end{description}

%%blk background.faults.03 from=716ec3fd47ac
Η συμπεριφορά που χωρίζει το βυζαντινό μοντέλο από τα υπόλοιπα έχει
όνομα. Είναι η \emph{αντιφατική δήλωση} (\en{equivocation}). Ο ίδιος
κόμβος λέει διαφορετικά πράγματα σε διαφορετικούς παραλήπτες, και οι
παραλήπτες δεν μπορούν να το διαπιστώσουν μόνοι τους. Το
Σχήμα~\ref{fig:equivocation} το δείχνει.

%%blk background.faults.fig.equivocation from=c119ec21a4c3
\dgfigure{fig:equivocation}%
{Κατάρρευση έναντι αντιφατικής δήλωσης: η μία συμπεριφορά που κάνει το
βυζαντινό μοντέλο ακριβό.}%
{Κατάρρευση έναντι αντιφατικής δήλωσης: η μία συμπεριφορά που κάνει το
βυζαντινό μοντέλο ακριβό. Αριστερά, ο $K_2$ σταματά. Οι $K_1$ και $K_3$
έχουν λιγότερη πληροφορία, αλλά όχι αντιφατική. Δεξιά, ο $K_2$ στέλνει
διαφορετική τιμή σε καθέναν και συνεχίζει να λειτουργεί. Ούτε ο $K_1$
ούτε ο $K_3$ μπορεί να το διαπιστώσει μόνος του. Το έμπιστο υλικό
μειώνει το φράγμα από $3f{+}1$ σε $2f{+}1$ ακριβώς επειδή εμποδίζει αυτή
τη συμπεριφορά. Η ταξινόμηση μεταφέρεται από την ανασκόπηση
\cite{amiri2024bedrock} (Παράρτ.~A, O7). Η πρωτογενής πηγή του φράγματος
είναι η γραμμή έργων έμπιστου υλικού στην οποία παραπέμπει η ανασκόπηση:
το TrInc, που το κατονομάζει, και άλλα στα οποία παραπέμπει με αριθμό
αναφοράς χωρίς να τα ονομάζει. Η εργασία \emph{δεν} τα έχει ελέγξει
ανεξάρτητα.}%
{% fig-equivocation — crash versus equivocation. Structure figure: dgflow.
% Left panel is the crash model, so every mark in it is dgCFT; right panel is
% the Byzantine model, so every mark in it is dgBFT. K1/K3 are ordinary nodes
% in both panels and stay dgbox grey. The two panels are 6 cm apart so that
% the pair comes out at \dgwidth.
\begin{dgflow}
  % ---- left: crash ------------------------------------------------------
  \node[text=dgCFT] at (1.5,3.35) {\textbf{\fl{equivocation.crash.title}}};
  \node[dgboxkey] (c2) at (0.0,1.2) {$K_2$};
  \node[dgbox]    (c1) at (3.0,2.4) {$K_1$};
  \node[dgbox]    (c3) at (3.0,0.0) {$K_3$};
  \draw[dgarrowkey] (c2) -- (c1)
        node[dgtag, fill=dgPaper, sloped, above, pos=0.5]
        {\fl{quote.open}$v=5$\fl{quote.close}};
  \draw[dgarrowkey] (c2) -- (c3)
        node[dgtag, fill=dgPaper, sloped, below, pos=0.5]
        {\fl{quote.open}$v=5$\fl{quote.close}};
  % K2 is gone: the cross sits just under the box it kills.
  \draw[dgarrowkey,-] (-0.30,0.65) -- (0.30,0.05);
  \draw[dgarrowkey,-] (-0.30,0.05) -- (0.30,0.65);
  \node[dgnote, anchor=north, align=center] at (1.5,-0.55)
       {\fl{equivocation.crash.note}};

  % ---- right: Byzantine -------------------------------------------------
  \node[text=dgBFT] at (7.5,3.35) {\textbf{\fl{equivocation.byzantine.title}}};
  \node[dgboxbad] (b2) at (6.0,1.2) {$K_2$};
  \node[dgbox]    (b1) at (9.0,2.4) {$K_1$};
  \node[dgbox]    (b3) at (9.0,0.0) {$K_3$};
  \draw[dgarrowbad] (b2) -- (b1)
        node[dgtag, fill=dgPaper, sloped, above, pos=0.5]
        {\fl{quote.open}$v=5$\fl{quote.close}};
  \draw[dgarrowbad] (b2) -- (b3)
        node[dgtag, fill=dgPaper, sloped, below, pos=0.5]
        {\fl{quote.open}$v=7$\fl{quote.close}};
  \node[dgnote, anchor=north, align=center] at (7.5,-0.55)
       {\fl{equivocation.byzantine.note}};
\end{dgflow}
}

%%blk background.faults.04 from=c7b1a25b883a
Ο Πίνακας~\ref{tab:dls} συνδυάζει τους δύο άξονες. Είναι το πιο
συμπυκνωμένο αποτέλεσμα του κεφαλαίου, και εξηγεί γιατί η μερική
συγχρονία δεν είναι δωρεάν.

%%blk background.faults.tab.dls from=e7e87ea4cae7
\begin{table}[t]
\centering
\caption[Το ελάχιστο πλήθος κόμβων ανά μοντέλο: το ελάχιστο $N$ για το
οποίο υπάρχει πρωτόκολλο συναίνεσης ανθεκτικό σε $t$ σφάλματα.]%
{Το ελάχιστο πλήθος κόμβων ανά μοντέλο: το ελάχιστο $N$ για το οποίο
υπάρχει πρωτόκολλο συναίνεσης ανθεκτικό σε $t$ σφάλματα. Ο πίνακας
αναπαράγει \textbf{τρεις από τις πέντε στήλες} του Πίνακα~I του
\cite{dwork1988consensus}. Η επιλογή αιτιολογείται στη σημείωση. Το
σύμβολο $\infty$ σημαίνει ότι δεν υπάρχει τέτοιο $N$, δηλαδή το FLP.
\emph{Δύο αναγνώσεις που αξίζουν:} στο σύγχρονο μοντέλο με
αυθεντικοποίηση ανέχεσαι \emph{οποιοδήποτε} πλήθος σφαλμάτων, και η
μερική συγχρονία ανεβάζει το κόστος του \en{fail-stop} από $t$ σε
$2t{+}1$, δηλαδή πληρώνεις για το ότι δεν ξέρεις πότε αρχίζει η
σταθερότητα.}
\label{tab:dls}
\small
\begin{tabular}{@{}l c c c@{}}
\toprule
\textbf{Είδος σφάλματος} & \textbf{Σύγχρονο} & \textbf{Ασύγχρονο} &
\textbf{Μερικώς σύγχρονο} \\
\midrule
\en{Fail-stop} (με την έννοια της πηγής) & $t$ & $\infty$ & $2t+1$ \\
Παραλείψεις               & $t$      & $\infty$ & $2t+1$ \\
Βυζαντινό με αυθεντικοποίηση & $t$   & $\infty$ & $3t+1$ \\
Βυζαντινό                 & $3t+1$   & $\infty$ & $3t+1$ \\
\bottomrule
\end{tabular}

\vspace{4pt}
\begin{minipage}{0.94\linewidth}\setstretch{1.08}\footnotesize
\textbf{Σημείωση ακρίβειας.} Η στήλη «Μερικώς σύγχρονο» αντιστοιχεί στην
περίπτωση όπου \emph{και} η επικοινωνία \emph{και} οι επεξεργαστές είναι
μερικώς σύγχρονοι. Το πρωτότυπο δίνει πέντε στήλες, επειδή διακρίνει και
τους συνδυασμούς όπου η μία διάσταση είναι σύγχρονη και η άλλη όχι. Η
έκδοση εδώ κρατά τις τρεις που χρησιμοποιεί η εργασία. Το ίδιο το άρθρο
βρίσκεται εκτός \en{corpus}. Το \emph{αντίτυπο αναφοράς} είναι το
ανοικτό θεσμικό αντίτυπο της ομάδας της Lynch στο MIT
(\texttt{groups.csail.mit.edu/tds/papers/Lynch/jacm88.pdf}). Από εκεί
διαβάσαμε τα μεγέθη του πίνακα, τους ορισμούς των $\Delta$ και $\Phi$,
τις δύο μορφές μερικής συγχρονίας και τον όρο GST. Δεν παρατίθεται
τίποτε αυτούσιο.

\smallskip
\textbf{Μία λέξη, δύο σημασίες.} Η πηγή ονομάζει την πρώτη γραμμή
\emph{\en{fail-stop}} και την ορίζει ως επεξεργαστή που εκτελεί σωστά,
μπορεί να σταματήσει οποτεδήποτε, και αφού σταματήσει δεν ξαναξεκινά.
\textbf{Δεν υπάρχει ρήτρα ανίχνευσης}. Δεν είναι αυτή η σημασία με την
οποία χρησιμοποιεί τη λέξη η Ενότητα~\ref{sec:faults} του κεφαλαίου,
ακολουθώντας το \cite{schneider1990implementing}, όπου το
\en{fail-stop} προσθέτει ότι οι υπόλοιποι κόμβοι αντιλαμβάνονται
\emph{αξιόπιστα} τη διακοπή. Η διαφορά κρίνει την ασύγχρονη στήλη: η
αξιόπιστη ανίχνευση είναι γνησίως ισχυρότερη υπόθεση από εκείνη που
κάνει ο πίνακας, και το $\infty$ της στήλης ανήκει στην ασθενέστερη
υπόθεση της πηγής. Η γραμμή επομένως φέρει το όνομα που της δίνει η
πηγή, και ο αναγνώστης πληροφορείται ποια από τις δύο σημασίες ισχύει.
Είναι η ίδια παγίδα με την Ενότητα~\ref{sec:smr}: μία λέξη, δύο
μετρήσεις.
\end{minipage}
\end{table}

%%blk background.smr.head from=30f775134ee8
\section{Συναίνεση, ατομική εκπομπή και SMR δεν είναι το ίδιο}
\label{sec:smr}

%%blk background.smr.01 from=3c65c3ca4e01
Τα τρία ονόματα χρησιμοποιούνται συχνά εναλλάξ. Δεν είναι συνώνυμα, και
η σύγχυση παράγει λάθος συγκρίσεις στα επόμενα κεφάλαια.

%%blk background.smr.list from=3299af9be4eb
\begin{description}[leftmargin=2.2em,style=nextline,font=\normalfont\bfseries]
  \item[Συναίνεση (\en{consensus}).] Οι κόμβοι συμφωνούν σε \emph{μία}
  τιμή.
  \item[Ατομική εκπομπή (\en{atomic broadcast}).] Οι κόμβοι παραδίδουν
  την \emph{ίδια ακολουθία} μηνυμάτων με την ίδια σειρά.
  \item[Αναπαραγωγή μηχανής καταστάσεων (SMR).] Η ατομική εκπομπή
  εφαρμόζεται σε ντετερμινιστική μηχανή, ώστε όλα τα αντίγραφα να
  καταλήξουν στην ίδια κατάσταση \cite{schneider1990implementing}.
\end{description}

%%blk background.smr.02 from=8216bb1b38e6
Το Σχήμα~\ref{fig:smr} συνοψίζει τη στοίβα. Η SMR είναι το επίπεδο στο
οποίο ζει η εργασία. Ο Schneider αναλύει την απαίτησή της σε δύο
συνθήκες. Κάθε μη ελαττωματικό αντίγραφο λαμβάνει κάθε αίτημα, και κάθε
μη ελαττωματικό αντίγραφο τα επεξεργάζεται με την ίδια σχετική σειρά. Η
πρώτη συνθήκη λέγεται \emph{συμφωνία}, η δεύτερη \emph{διάταξη}.
Πρόσεξε ότι η δεύτερη μιλά για \emph{σχετική} σειρά, και όχι για
καθολική χρονοσφραγίδα.

%%blk background.smr.03 from=4501fac53591
Πάνω από όλα αυτά στέκει μια υπόθεση που σπάει εύκολα στην πράξη. Η
ακολουθία των αιτημάτων πρέπει να καθορίζει πλήρως τις εξόδους της
μηχανής, και τίποτε άλλο δεν επιτρέπεται να τις επηρεάζει
\cite{schneider1990implementing}. Αν η μηχανή καλέσει γεννήτρια τυχαίων
αριθμών ή ρολόι, η SMR παύει να ισχύει.

%%blk background.smr.fig.smr from=8749db49f0a2
\dgfigure{fig:smr}%
{Η στοίβα: πού ζει το καθένα.}%
{Η στοίβα: πού ζει το καθένα. Η συναίνεση λύνει τη σειρά. Η μηχανή
καταστάσεων εκτελεί. Όλα τα πρωτόκολλα της εργασίας κατοικούν στο μεσαίο
κουτί, και γι' αυτό συγκρίνονται. Η ντετερμινιστικότητα των μηχανών
είναι υπόθεση και όχι συμπέρασμα \cite{schneider1990implementing}.}%
{% fig-smr — where each layer lives. Structure figure: dgflow.
% The middle box is the subject of the whole thesis, so it is dgboxkey; the
% clients and the replicas are ordinary components, so they are dgbox. Both
% CFT and BFT protocols live in the middle box, which is why it cannot carry
% either series colour.
\begin{dgflow}
  \node[dgbox]    (cl) at (0.0,0.0) {\fl{smr.clients}};
  \node[dgboxkey] (co) at (4.0,0.0)
       {\bfseries\fl{smr.consensus}\\\mdseries\fl{smr.agreement.order}};
  \node[dgbox]    (sm1) at (8.4, 0.95) {\fl{smr.machine} 1};
  \node[dgbox]    (sm2) at (8.4, 0.00) {\fl{smr.machine} 2};
  \node[dgbox]    (sm3) at (8.4,-0.95) {\fl{smr.machine} 3};
  \draw[dgarrow] (cl) -- (co);
  \draw[dgarrow] (co) -- (sm1);
  \draw[dgarrow] (co) -- (sm2);
  \draw[dgarrow] (co) -- (sm3);
  \node[dgnote, anchor=north, align=center] at (8.4,-1.45)
       {\fl{smr.deterministic}\\\fl{smr.same.io}};
  \node[dgnote, anchor=north] at (4.0,-1.15)
       {\fl{smr.here.live} Paxos, Raft, PBFT, HotStuff};
\end{dgflow}
}

%%blk background.smr.sub.trap from=d55660872921
\subsection{Μια παγίδα αρίθμησης: $2t{+}1$ και $3t{+}1$ δεν μετρούν το
ίδιο}

%%blk background.smr.04 from=7b142cd8d1bf
Ο Schneider γράφει ότι με βυζαντινά σφάλματα χρειάζονται τουλάχιστον
$2t{+}1$ αντίγραφα, ενώ με σφάλματα \en{fail-stop} αρκούν $t{+}1$
\cite{schneider1990implementing}. Η δεύτερη τιμή προϋποθέτει ότι η πτώση
είναι \emph{ανιχνεύσιμη}. Τότε αρκεί να πιστέψεις οποιονδήποτε ζωντανό
κόμβο. Το Κεφάλαιο~\ref{ch:intro} γράφει $3f{+}1$ για το βυζαντινό
μοντέλο. Τα δύο δεν συγκρούονται, επειδή μετρούν διαφορετικά πράγματα.

%%blk background.smr.trap.list from=51d196acccb0
\begin{itemize}
  \item Ο Schneider μετρά αντίγραφα των οποίων οι \emph{έξοδοι
  ψηφίζονται}. Με $2t{+}1$ αντίγραφα, η πλειοψηφία των εξόδων μένει
  σωστή ακόμη και μετά από $t$ σφάλματα. Η συμφωνία και η διάταξη
  θεωρούνται ήδη λυμένες.
  \item Το $3f{+}1$ είναι το φράγμα για να \emph{λυθεί} η συμφωνία,
  δηλαδή το φράγμα για το ίδιο το πρωτόκολλο συναίνεσης
  \cite{pease1980reaching}.
\end{itemize}

%%blk background.smr.05 from=ebe438b5427c
Η σύγχυση των δύο είναι εύκολη, επειδή και τα δύο μεγέθη γράφονται ως
συνάρτηση του πλήθους των σφαλμάτων. Η εργασία χρησιμοποιεί το $3f{+}1$
παντού, επειδή μετρά πρωτόκολλα συναίνεσης και όχι ψηφοφορία εξόδων.

%%blk background.bounds.head from=8abb95caf712
\section{Τι είναι αδύνατο και τι κοστίζει}
\label{sec:bounds}

%%blk background.bounds.roadmap from=dd5e7b13c920
Τέσσερα αποτελέσματα οριοθετούν κάθε πρωτόκολλο της εργασίας. Το πρώτο
λέει ότι σε πλήρως ασύγχρονο σύστημα κανένα ντετερμινιστικό πρωτόκολλο
δεν τερματίζει πάντοτε. Το δεύτερο δίνει το ελάχιστο πλήθος κόμβων που
ανέχεται δεδομένο πλήθος βυζαντινών σφαλμάτων. Το τρίτο δίνει το κόστος
σε μηνύματα. Το τέταρτο δίνει το ελάχιστο πλήθος βημάτων επικοινωνίας.
Παρουσιάζονται με αυτή τη σειρά, γιατί το καθένα στενεύει τον χώρο που
άφησε ανοιχτό το προηγούμενο.

%%blk background.bounds.sub.flp from=3e265ae78fed
\subsection{FLP}

%%blk background.bounds.01 from=734383c3efc2
Το θεμελιώδες αρνητικό αποτέλεσμα διατυπώνεται σε μία γραμμή. Οι
Fischer, Lynch και Paterson αποδεικνύουν ότι «\en{No consensus protocol
is totally correct in spite of one fault}» \cite{fischer1985impossibility}.
Το πρωτότυπο βρίσκεται στο \en{corpus} και η φράση παρατίθεται αυτούσια.

%%blk background.bounds.02 from=ed658280bad5
Η μία γραμμή όμως κουβαλά δύο υποθέσεις, και το ίδιο το άρθρο τις
δηλώνει ρητά. Το σύστημα είναι \emph{πλήρως ασύγχρονο} και το πρωτόκολλο
\emph{ντετερμινιστικό}. Η δεύτερη υπόθεση είναι η πόρτα εξόδου. Τα
τυχαιοποιημένα πρωτόκολλα δεν εμπίπτουν στο θεώρημα, και το ίδιο το
άρθρο παραπέμπει σε τέτοια εργασία. Όλα τα πρωτόκολλα της εργασίας είναι
ντετερμινιστικά, άρα το θεώρημα τα δεσμεύει.

%%blk background.bounds.03 from=36f70b480c6b
Τρεις ακρίβειες αλλάζουν τη διατύπωση του αποτελέσματος. Και τις τρεις
τις διαβάσαμε στο ίδιο το άρθρο.

%%blk background.bounds.flp.list from=a7b9801c9c72
\begin{enumerate}
  \item Το άρθρο \textbf{δεν} ορίζει την ορθότητα μέσω των επώνυμων
  ιδιοτήτων «συμφωνία» και «εγκυρότητα». Η λέξη \en{validity} δεν
  εμφανίζεται καθόλου, ενώ η \en{agreement} εμφανίζεται μόνο ως όνομα
  του προβλήματος. Ορίζει ένα πρωτόκολλο ως \emph{μερικώς ορθό} με δύο
  συνθήκες. Καμία προσβάσιμη διαμόρφωση δεν έχει πάνω από μία τιμή
  απόφασης, και για κάθε τιμή υπάρχει προσβάσιμη διαμόρφωση που την
  αποφασίζει. Η δεύτερη συνθήκη είναι η \emph{μη τετριμμένη}, που είναι
  ασθενέστερη από την εγκυρότητα.
  \item Το άρθρο \textbf{δεν} εξετάζει βυζαντινά σφάλματα, και το
  δηλώνει ρητά. Υποθέτει αξιόπιστο σύστημα μηνυμάτων και μία μόνο
  κατάρρευση. Αυτό κάνει το αποτέλεσμα \emph{ισχυρότερο}. Αν η συναίνεση
  είναι αδύνατη στο ασθενέστερο μοντέλο σφαλμάτων, είναι κατά μείζονα
  λόγο αδύνατη στο βυζαντινό.
  \item Η απαίτηση τερματισμού είναι \textbf{εξαιρετικά ασθενής}, και το
  άρθρο το δηλώνει. Αρκεί \emph{κάποιος} κόμβος να αποφασίσει τελικά,
  και όχι όλοι. Το αρνητικό αποτέλεσμα δεν οφείλεται λοιπόν ούτε σε
  κακόβουλους κόμβους ούτε σε χαμένα μηνύματα. Οφείλεται μόνο στο ότι η
  σιωπή δεν ξεχωρίζει από την αργοπορία.
\end{enumerate}

%%blk background.bounds.sub.third from=b7ca28913a8a
\subsection{Το φράγμα του ενός τρίτου}

%%blk background.bounds.04 from=0fc36e0f0a99
Το φράγμα $n \ge 3f+1$ για βυζαντινή συμφωνία χωρίς αυθεντικοποίηση
δημοσιεύεται πρώτη φορά το 1980, με τη μορφή του προβλήματος της
\emph{διαδραστικής συνέπειας} \cite{pease1980reaching}. Το όνομα λέει
την ίδια συνθήκη από την ανάποδη. Οι ελαττωματικοί πρέπει να είναι
λιγότεροι από το ένα τρίτο των κόμβων, άρα οι σωστοί ξεπερνούν τα δύο
τρίτα. Δύο χρόνια αργότερα οι ίδιοι συγγραφείς το ξαναδιατυπώνουν ως
αλληγορία των βυζαντινών στρατηγών \cite{lamport1982byzantine}.

%%blk background.bounds.05 from=81bbc88f29fa
Το ίδιο άρθρο του 1980 εξετάζει ήδη και την \emph{αυθεντικοποιημένη}
περίπτωση \cite{pease1980reaching}. Εκεί οι ελαττωματικοί κόμβοι μπορούν
να σιωπούν, αλλά δεν μπορούν να παραποιούν όσα αναμεταδίδουν. Το φράγμα
του ενός τρίτου δεν ισχύει σε αυτή την περίπτωση.

%%blk background.bounds.06 from=23bf0896781f
Η αλληγορία του 1982 προσθέτει την καθαρότερη διατύπωση. Ο αλγόριθμος
SM($m$) λύνει το πρόβλημα για οποιοδήποτε πλήθος στρατηγών, αρκεί η
υπογραφή ενός πιστού στρατηγού να μην πλαστογραφείται
\cite{lamport1982byzantine}. Το ίδιο άρθρο βάζει μια ρήτρα. Το πρόβλημα
είναι κενό με λιγότερους από $m{+}2$ στρατηγούς. Το αποτέλεσμα των
υπογραφών δεν είναι λοιπόν συνεισφορά του 1982, αλλά του 1980.

%%blk background.bounds.07 from=24627e844d97
Το φράγμα $3f{+}1$ που χρησιμοποιούν όλα τα πρωτόκολλα της εργασίας δεν
προέρχεται λοιπόν από την απουσία υπογραφών. Προέρχεται από τη
\emph{μερική συγχρονία}, όπως δείχνει η γραμμή «Βυζαντινό με
αυθεντικοποίηση» του Πίνακα~\ref{tab:dls}. Εκεί η απαίτηση είναι
$n \ge 3t{+}1$ ακόμη και όταν τα μηνύματα είναι υπογεγραμμένα.

%%blk background.bounds.quote from=d8bd4398be9d
\begin{quote}\small
Αν σε ρωτήσουν «γιατί οι υπογραφές δεν σε γλιτώνουν από το $3f{+}1$;», η
απάντηση είναι αυτή η γραμμή του πίνακα. Οι υπογραφές αρκούν στο
\emph{σύγχρονο} μοντέλο. Στο μερικώς σύγχρονο δεν αρκούν.
\end{quote}

%%blk background.bounds.sub.messages from=14ef45ab1fd7
\subsection{Το κόστος σε μηνύματα}

%%blk background.bounds.08 from=6a2ef96a77cd
Η αδυναμία δεν είναι το μόνο φράγμα. Υπάρχει και κάτω φράγμα στο πλήθος
των μηνυμάτων. Οι Dolev και Reischuk αποδεικνύουν ότι κάθε
\textbf{ντετερμινιστικός} αλγόριθμος βυζαντινής συμφωνίας που ανέχεται
$t$ σφάλματα έχει εκτέλεση στην οποία οι σωστοί κόμβοι στέλνουν
τετραγωνικό ως προς $t$ πλήθος μηνυμάτων \cite{dolev1985bounds}.

%%blk background.bounds.09 from=d5b89e6bdb84
Τρία σημεία απαιτούν προσοχή, και το κεφάλαιο δηλώνει και τα τρία εδώ
αντί να τα αποσιωπήσει.

%%blk background.bounds.messages.list from=503bf627cf56
\begin{itemize}
  \item Το θεώρημα \textbf{δεν} περιορίζεται στους μη
  αυθεντικοποιημένους αλγορίθμους. Η υπόθεσή του δεν περιέχει ρήτρα
  αυθεντικοποίησης, άρα καλύπτει και τις δύο περιπτώσεις. Πολλές
  δευτερογενείς πηγές το γράφουν λανθασμένα.
  \item Ο περιορισμός στον \textbf{ντετερμινισμό} δεν είναι
  διακοσμητικός. Τα τυχαιοποιημένα πρωτόκολλα μπορούν να πετύχουν
  υποτετραγωνική επικοινωνία, και πληρώνουν με πιθανότητα σφάλματος.
  Αυτό δεν ισχύει πάντα. Απέναντι σε \emph{ισχυρά προσαρμοστικό}
  αντίπαλο, που διαγράφει μηνύματα εκ των υστέρων, το τετραγωνικό φράγμα
  ισχύει και για αυτά \cite{abraham2019ccba}. Δεν σπάνε λοιπόν το
  φράγμα. Βγαίνουν από την υπόθεσή του, και μόνο υπό ασθενέστερο μοντέλο
  αντιπάλου. Όλα τα πρωτόκολλα της εργασίας είναι ντετερμινιστικά.
  \item Το ίδιο το άρθρο του περιοδικού βρίσκεται επί πληρωμή και εκτός
  \en{corpus}. Ελέγξαμε το περιεχόμενο στην προγενέστερη εκδοχή
  συνεδρίου, και η αρίθμηση των θεωρημάτων μπορεί να διαφέρει.
\end{itemize}

%%blk background.bounds.10 from=173331c5a7cb
Η συνέπεια αφορά άμεσα τα κεφάλαια των βυζαντινών πρωτοκόλλων. Όταν το
HotStuff διαφημίζει \emph{γραμμική} αλλαγή όψης, ο ισχυρισμός ισχύει και
για την αλλαγή όψης, εφόσον ο ηγέτης που αναλαμβάνει είναι σωστός και το
δίκτυο σταθερό \cite{yin2019hotstuff}. Το κάτω φράγμα δεν παραβιάζεται.
Μετατοπίζεται στη χειρότερη περίπτωση. Το ίδιο άρθρο αναφέρει $O(fn)$
\emph{αυθεντικοποιητές} πάνω σε $f$ διαδοχικούς ελαττωματικούς ηγέτες,
και αυτό είναι $O(n^{2})$ όταν το $f$ μεγαλώνει μαζί με το $n$.

%%blk background.bounds.11 from=4db334762ea8
\textbf{Προσοχή στη μονάδα.} Το HotStuff μετρά \en{authenticators}, ενώ
το φράγμα των Dolev και Reischuk μετρά \emph{μηνύματα}. Τα δύο μεγέθη
δεν ταυτίζονται, και το ίδιο το άρθρο του HotStuff το δηλώνει. Ένα
μήνυμα που κουβαλά $n$ υπογραφές μετρά όσο $n$ μηνύματα με μία υπογραφή
το καθένα. Η αντιπαράθεση εδώ είναι ποιοτική και όχι λογιστική.

%%blk background.bounds.sub.twostep from=4b21385912d2
\subsection{Το φράγμα των δύο βημάτων}

%%blk background.bounds.12 from=132bac2fe5d2
Η τελευταία κατηγορία φραγμάτων αφορά τον αριθμό των φάσεων. \emph{Δύο
βήματα} σημαίνει ότι μια εντολή δεσμεύεται με δύο μόνο διαδοχικές
φάσεις. Αυτές είναι η πρόταση του ηγέτη και ένας γύρος ψήφων.

%%blk background.bounds.13 from=8b82c5ddf6c7
\textbf{Προσοχή στη σύμβαση αρίθμησης.} Η ανασκόπηση
\cite{amiri2024bedrock} μετρά το PBFT με \emph{τρεις} φάσεις διάταξης
και το FaB με \emph{δύο}. Ο Πίνακας~\ref{tab:delays} ακολουθεί τη
σύμβαση του \cite{malkhi2023hotstuff2}\prep, όπου το PBFT εμφανίζεται με
δύο. Οι δύο αριθμήσεις δεν συγκρίνονται μεταξύ τους.

%%blk background.bounds.14 from=a21d7995b2e7
Η εξοικονόμηση της μίας φάσης κόστιζε αρχικά $5f{+}1$ κόμβους
\cite{martin2006fab}\unver. Με αυθεντικοποίηση αρκούν $5f{-}1$, και αυτό είναι
το κάτω φράγμα \cite{kuznetsov2021optimal,abraham2021goodcase}\unver.

%%blk background.bounds.15 from=a1c8a7aed969
Η ανασκόπηση \cite{amiri2024bedrock} μεταφέρει και τα δύο μεγέθη, και
από εκεί τα πήρε η εργασία. Το κείμενο κατονομάζει τις πρωτογενείς
πηγές, αλλά η εργασία δεν τις έχει ελέγξει ανεξάρτητα. Ο
Πίνακας~\ref{tab:bounds} το δηλώνει.

%%blk background.bounds.tab.bounds from=db1b365d08c0
\begin{table}[t]
\centering
\caption[Σύνοψη των φραγμάτων που χρησιμοποιεί η εργασία, με την
πρωτογενή τους πηγή.]%
{Σύνοψη των φραγμάτων που χρησιμοποιεί η εργασία, με την πρωτογενή τους
πηγή. Η στήλη «Τεκμηρίωση» είναι μέρος του τεκμηρίου και έχει τρεις
τιμές. \emph{Corpus}: το άρθρο υπάρχει και το διαβάσαμε.
\emph{Ευρετήριο}: τα βιβλιογραφικά επαληθεύτηκαν σε DOI ή σε ευρετήριο,
και το περιεχόμενο σε ανοικτό θεσμικό ή αυθεντικό αντίτυπο του
συγγραφέα, ή σε προγενέστερη εκδοχή. \emph{Μη επαληθευμένη}: τα στοιχεία
αντλήθηκαν από τον κατάλογο αναφορών του \cite{amiri2024bedrock}. Όπου η
στήλη δεν λέει \emph{corpus}, το κείμενο δεν παραθέτει τίποτε
αυτούσιο.}
\label{tab:bounds}
\small
\begin{tabular}{@{}l l c@{}}
\toprule
\textbf{Φράγμα} & \textbf{Πρωτογενής πηγή} & \textbf{Τεκμηρίωση} \\
\midrule
Αδυναμία σε ασύγχρονο, 1 κατάρρευση & \cite{fischer1985impossibility} & corpus \\
$n \ge 3f+1$, βυζαντινό & \cite{pease1980reaching} & ευρετήριο \\
Υπογραφές: το φράγμα καταρρέει (σύγχρονο) &
  \cite{pease1980reaching,lamport1982byzantine} & ευρετήριο \\
$n \ge 2f+1$ / $3f+1$ σε μερική συγχρονία & \cite{dwork1988consensus} & ευρετήριο \\
$\Omega(t^2)$ μηνύματα, ντετερμινιστικά & \cite{dolev1985bounds} & ευρετήριο \\
$5f{+}1$ για μία φάση λιγότερη & \cite{martin2006fab}\unver\ & μη επαληθευμένη \\
$5f{-}1$ κάτω φράγμα δύο βημάτων &
  \cite{kuznetsov2021optimal,abraham2021goodcase}\unver\ & μη επαληθευμένη \\
\bottomrule
\end{tabular}
\end{table}

%%blk background.bounds.16 from=0e496d95948a
Ο Πίνακας~\ref{tab:bounds} συγκεντρώνει τα φράγματα του κεφαλαίου μαζί
με την τεκμηρίωσή τους.

%%blk background.quorums.head from=8808e23da0f0
\section{Απαρτίες: η μία ιδιότητα που κάνει τη δουλειά}
\label{sec:quorums}

%%blk background.quorums.01 from=bb323c478274
Μια \emph{απαρτία} (\en{quorum}) είναι το ελάχιστο σύνολο κόμβων που
κάνει μια απόφαση οριστική. Η μοναδική ιδιότητα που χρειάζεται είναι η
τομή.

%%blk background.quorums.sub.derivation from=6a02fae17460
\subsection{Η παραγωγή}

%%blk background.quorums.02 from=f2d0985c53ca
Το Σχήμα~\ref{fig:quorum} του Κεφαλαίου~\ref{ch:intro} δίνει τα μεγέθη.
Εδώ είναι το επιχείρημα από το οποίο προκύπτουν, σε τρία βήματα.

%%blk background.quorums.list from=888b33260d40
\begin{enumerate}
  \item Δεν μπορείς να περιμένεις πάνω από $n-f$ απαντήσεις. Οι $f$
  ελαττωματικοί μπορεί να σιωπούν, και η σιωπή δεν ξεχωρίζει από την
  αργοπορία. Άρα η απαρτία έχει μέγεθος $q = n-f$.
  \item Δύο απαρτίες τέμνονται σε τουλάχιστον $2q - n = n - 2f$ κόμβους.
  \item Στο βυζαντινό μοντέλο η τομή πρέπει να περιέχει τουλάχιστον έναν
  \emph{τίμιο} κόμβο. Στην τομή κρύβονται έως $f$ βυζαντινοί, άρα
  $n - 2f > f$, δηλαδή $n \ge 3f+1$.
\end{enumerate}

%%blk background.quorums.03 from=4d420e6f4761
Στο μοντέλο κατάρρευσης το βήμα~3 δεν υπάρχει. Αρκεί $n - 2f > 0$,
δηλαδή $n \ge 2f+1$. Αυτό το ένα βήμα είναι ολόκληρη η διαφορά κόστους
μεταξύ των δύο μοντέλων.

%%blk background.quorums.quote from=6436242bcf09
\begin{quote}\small
\textbf{Τι είναι και τι δεν είναι αυτό.} Η παραπάνω παραγωγή είναι
επιχείρημα \emph{σχεδίασης απαρτίας}. Δείχνει ότι οι απαρτίες μεγέθους
$n-f$ χρειάζονται $n \ge 3f+1$. Δεν είναι η γενική απόδειξη
αδυνατότητας, που ανήκει στο \cite{pease1980reaching}.
\end{quote}

%%blk background.quorums.sub.intersection from=34a4b93980f6
\subsection{Η τομή απαιτείται λιγότερο απ' όσο νομίζαμε}

%%blk background.quorums.04 from=e1b30045befd
Από το Paxos και έπειτα, η πλειοψηφική απαρτία θεωρούνταν αναγκαία επί
Οι Howard, Malkhi και Spiegelman δείχνουν ότι
δεν είναι \cite{howard2016flexiblepaxos}. Στο Paxos, η τομή απαιτείται
μόνο \emph{μεταξύ} των δύο φάσεων. Οι απαρτίες της Φάσης~1 δεν
χρειάζεται να τέμνονται μεταξύ τους, ούτε εκείνες της Φάσης~2 μεταξύ
τους.

%%blk background.quorums.05 from=38f7d807d120
Η πρακτική συνέπεια είναι ότι μπορείς να μικρύνεις την απαρτία της
Φάσης~2, δηλαδή της σταθερής κατάστασης. Πληρώνεις με μεγαλύτερη απαρτία
στην εκλογή ηγέτη. Το κεφάλαιο του Paxos επανέρχεται σε αυτό, επειδή
αλλάζει την ανάγνωση της αριθμητικής του Σχήματος~\ref{fig:quorum}.

%%blk background.correctness.head from=7b9e028490ea
\section{Κριτήρια ορθότητας}
\label{sec:correctness}

%%blk background.correctness.roadmap from=1f66b7a0fe6c
Δύο ερωτήματα κρίνουν αν μια αναπαραγόμενη υπηρεσία είναι ορθή, και
είναι ξεχωριστά ερωτήματα. Το πρώτο είναι τι εγγυάται το πρωτόκολλο για
τη σειρά που χτίζει· αυτό είναι η ασφάλεια και η ζωντάνια. Το δεύτερο
είναι τι δικαιούται να παρατηρήσει ο πελάτης από έξω· αυτό είναι η
γραμμικοποιησιμότητα. Η ενότητα τα παίρνει με αυτή τη σειρά, γιατί το
δεύτερο διατυπώνεται με όρους του πρώτου.

%%blk background.correctness.sub.safety from=1273024c3aa1
\subsection{Ασφάλεια και ζωτικότητα}

%%blk background.correctness.01 from=fa7fec866fba
Ένα πρωτόκολλο κρίνεται από δύο κατηγορίες ιδιοτήτων. Η \emph{ασφάλεια}
λέει ότι κάτι κακό δεν συμβαίνει ποτέ. Η \emph{ζωτικότητα} λέει ότι κάτι
καλό συμβαίνει τελικά.

%%blk background.correctness.02 from=a8a76350284b
Η ασυμμετρία που διέπει όλα τα πρωτόκολλα της εργασίας είναι απλή. Η
ασφάλεια διατηρείται πάντοτε, ακόμη και σε πλήρη διαμέριση δικτύου. Η
ζωτικότητα θυσιάζεται. Η επιλογή είναι σχεδιαστική και όχι ατύχημα. Ένα
σύστημα που σταματά διορθώνεται, ενώ ένα σύστημα που έγραψε δύο
διαφορετικές τιμές δεν διορθώνεται.

%%blk background.correctness.sub.lin from=f31ac030bd03
\subsection{Γραμμικοποιησιμότητα}

%%blk background.correctness.03 from=4d58102b4893
Οι Herlihy και Wing ορίζουν το κριτήριο που περιμένει ο πελάτης
\cite{herlihy1990linearizability}. Κάθε πράξη φαίνεται να συμβαίνει
ακαριαία σε κάποιο σημείο ανάμεσα στην κλήση της και στην απόκρισή της.
Οι πράξεις που δεν επικαλύπτονται χρονικά διατηρούν τη σειρά τους.

%%blk background.correctness.04 from=6ee23d52022b
Το δεύτερο σκέλος κάνει όλη τη δουλειά. Είναι περιορισμός
\emph{πραγματικού χρόνου}, και γι' αυτό η γραμμικοποιησιμότητα κοστίζει
συντονισμό.

%%blk background.correctness.05 from=d395b473ce06
Δύο ιδιότητες τη χωρίζουν από τη σειριοποιησιμότητα των βάσεων
δεδομένων, και οι δύο έχουν σημασία για την εργασία.

%%blk background.correctness.list from=7d5c69fe512d
\begin{description}[leftmargin=2.2em,style=nextline,font=\normalfont\bfseries]
  \item[Τοπικότητα (\en{locality}).] Ένα σύστημα είναι γραμμικοποιήσιμο
  αν κάθε αντικείμενό του είναι. Η ιδιότητα \emph{συντίθεται}. Ούτε η
  σειριοποιησιμότητα ούτε η αυστηρή σειριοποιησιμότητα είναι τοπικές
  \cite{herlihy1990linearizability}.
  \item[Μη φραγή (\en{nonblocking}).] Μια εκκρεμής κλήση ολικής πράξης
  δεν υποχρεώνεται ποτέ να περιμένει άλλη εκκρεμή κλήση. Η
  σειριοποιησιμότητα είναι εγγενώς φράσσουσα
  \cite{herlihy1990linearizability}.
\end{description}

%%blk background.correctness.06 from=5c3e44100634
Οι ίδιοι συγγραφείς περιγράφουν τη γραμμικοποιησιμότητα ως ειδική
περίπτωση της αυστηρής σειριοποιησιμότητας, όπου κάθε δοσοληψία
περιορίζεται σε μία πράξη πάνω σε ένα αντικείμενο.

%%blk background.performance.head from=3c4350eec769
\section{Τι σημαίνει «επίδοση» για ένα πρωτόκολλο}
\label{sec:performance}

%%blk background.performance.01 from=7df6d5147d38
Αυτή η ενότητα είναι σύντομη και φέρει βάρος. Ορίζει, μία φορά, τη
διαφορά ανάμεσα σε δύο μεγέθη που το Κεφάλαιο~\ref{ch:intro} έδειξε ότι
συγχέονται.

%%blk background.performance.list from=c4e4829ca152
\begin{description}[leftmargin=2.2em,style=nextline,font=\normalfont\bfseries]
  \item[Φάσεις.] Πόσοι διαδοχικοί γύροι πρότασης και ψήφου χρειάζονται
  για να δεσμευτεί μια εντολή. Είναι ιδιότητα του \emph{πρωτοκόλλου} και
  \emph{συγκρίσιμη} μεταξύ άρθρων, εφόσον δηλωθεί η σύμβαση αρίθμησης.
  \item[Καθυστερήσεις μηνυμάτων.] Πόσες διαδοχικές αποστολές
  χρειάζονται για να δεσμευτεί μια εντολή. Στο HotStuff μία φάση
  αντιστοιχεί σε δύο καθυστερήσεις. Αλλού η αντιστοιχία διαφέρει, και
  γι' αυτό δηλώνεται.
  \item[Πραγματικός χρόνος.] Πόσα χιλιοστά του δευτερολέπτου
  μεσολαβούν. Εξαρτάται από την υλοποίηση, το υλικό και την τοπολογία
  του δικτύου. Δεν είναι συγκρίσιμος μεταξύ άρθρων. Η εργασία δίνει
  δικές της μετρήσεις πραγματικού χρόνου στο κεφάλαιο της αποτίμησης.
\end{description}

%%blk background.performance.02 from=77bc94736876
Ο Πίνακας~\ref{tab:delays} δίνει το πρώτο μέγεθος, σε \emph{φάσεις}, και
η σημείωσή του δηλώνει τη μετατροπή σε καθυστερήσεις. Είναι ο μόνος
πίνακας επίδοσης του κεφαλαίου, επειδή είναι ο μόνος που δεν εξαρτάται
από περιβάλλον.

%%blk background.performance.tab.delays from=03a44590972b
\begin{table}[t]
\centering
\caption[Φάσεις στην ευτυχή διαδρομή: πόσες φάσεις χρειάζονται για μία
απόφαση, με σωστό ηγέτη και μετά την GST.]%
{Φάσεις στην ευτυχή διαδρομή: πόσες φάσεις χρειάζονται για μία απόφαση,
με σωστό ηγέτη και μετά την GST. Δεν υπάρχουν εδώ χιλιοστά του
δευτερολέπτου, και αυτό είναι το νόημα της ενότητας. Πηγές:
\cite{malkhi2023hotstuff2}\prep\ (Πίνακας~1) για τα τέσσερα βυζαντινά,
και \cite{ongaro2014raft} \S5.3 για τη γραμμή Multi-Paxos / Raft, όπου
σε σταθερή κατάσταση «\en{a new entry can be replicated with a single
round of RPCs to a majority of the cluster}».}
\label{tab:delays}
\small
\begin{tabular}{@{}l c c@{}}
\toprule
\textbf{Πρωτόκολλο} & \textbf{Φάσεις} & \textbf{Μοντέλο} \\
\midrule
Multi-Paxos / Raft & 1 & κατάρρευση \\
PBFT               & 2 & βυζαντινό \\
Tendermint         & 2 & βυζαντινό, μη αποκρίσιμο \\
HotStuff           & 3 & βυζαντινό \\
HotStuff-2\prep    & 2 & βυζαντινό \\
\bottomrule
\end{tabular}

\vspace{4pt}
\begin{minipage}{0.94\linewidth}\setstretch{1.08}\footnotesize
\textbf{Προσοχή στη λογιστική.} Ο πίνακας μετρά \emph{φάσεις}, με την
ίδια έννοια και την ίδια πηγή που χρησιμοποιεί ο
Πίνακας~\ref{tab:viewchange} του Κεφαλαίου~\ref{ch:intro}. Στο HotStuff
κάθε φάση αντιστοιχεί σε \emph{δύο} διαδοχικές καθυστερήσεις μηνυμάτων,
δηλαδή πρόταση του ηγέτη και ψήφοι. Αυτό το δηλώνει η ίδια η πηγή
\emph{για το HotStuff}, άρα οι τρεις φάσεις του είναι έξι καθυστερήσεις.
Η αντιστοιχία \textbf{δεν} είναι καθολική. Το PBFT εμφανίζεται με $2$
επειδή έχει δύο γύρους ψηφοφορίας μετά την πρόταση, δηλαδή τρεις
καθυστερήσεις συνολικά και όχι τέσσερις. Ο πίνακας συγκρίνει λοιπόν
\emph{φάσεις όπως τις μετρά η πηγή}, και όχι καθυστερήσεις. Η γραμμή
Multi-Paxos / Raft δεν προέρχεται από την ίδια πηγή και σημειώνεται
χωριστά στη λεζάντα. Ο αριθμός των φάσεων \textbf{δεν} μεταφράζεται σε
χρόνο χωρίς να δηλωθεί το δίκτυο. Τρεις φάσεις σε LAN κοστίζουν λιγότερο
από δύο σε δίκτυο ευρείας περιοχής.
\end{minipage}
\end{table}

%%blk background.performance.03 from=e25074d85178
Η ενότητα κλείνει με τον κανόνα που θα εφαρμοστεί σε κάθε επόμενο
κεφάλαιο. Όταν συγκρίνονται πρωτόκολλα, συγκρίνονται φάσεις. Όταν
συγκρίνονται υλοποιήσεις, συγκρίνονται χρόνοι --- και μόνο εφόσον
δηλωθεί το περιβάλλον. Το πλήρες πλαίσιο εγκυρότητας, τα επίπεδα
συγκρισιμότητας και οι απειλές εγκυρότητας αναπτύσσονται στο κεφάλαιο
της μεθοδολογίας, όπως υπόσχεται το \S1.4.
